% Figure area_law_dyadic_layers.
% Owner: chapter/ch24_peps_area_law_dyadic_layers.tex (def:area_law_dyadic_layers).
% Used by: gapped-2d-notes AL Sec. 2 (dyadic layers); the note says that this D_k
% is a distance layer, not the exceptional set D.
% Source: 10-geometry.tex, lines 148--200 (N_k, D_k = N_{k+1} minus N_k);
% no source figure. Specialization drawn: Z is one endpoint and C_0 = 1.
% Ink-to-index: every dot is one half-open scale-k cell C_{k,z}; the red dot
% is the occupied index c_k(z) in Z_k. The pink enclosure is the 3-by-3
% block P_k, so its union is N_k. Each dashed square is one parent cell
% C_{k+1,p} of P_{k+1}, i.e. its four scale-k children; the nine of them
% make N_{k+1}, marked by the bracket. The tinted ring is the layer D_k,
% the union of the cells in I_k = Ch(P_{k+1}) minus P_k. N_k is not aligned
% with the parent squares, which is why D_k cuts through three of them.
% No strings, legs or contractions: this is a geometry picture with empty
% boundary signature. Not to scale; the plane continues past the frame.
\begin{center}
    \begin{tenkz}[size=s, rows={ket,ket,ket,ket,ket,ket,ket,ket}, cols=8, bonds=none]
        \tndeclare{species}{endpoint}{hue=source:red}
        \tn[at=(5,5), skin=dot, species=endpoint]{}
        \tnmark[form=enclosure, species=secondary, tint]{(2,2) .. (7,7) - (4,4) .. (6,6)}{}
        \tnmark[form=enclosure, species=selected, tint]{(4,4) .. (6,6)}{}
        \tnmark[form=enclosure, species=complement]{(2,2) .. (3,3)}{}
        \tnmark[form=enclosure, species=complement]{(2,4) .. (3,5)}{}
        \tnmark[form=enclosure, species=complement]{(2,6) .. (3,7)}{}
        \tnmark[form=enclosure, species=complement]{(4,2) .. (5,3)}{}
        \tnmark[form=enclosure, species=complement]{(4,4) .. (5,5)}{}
        \tnmark[form=enclosure, species=complement]{(4,6) .. (5,7)}{}
        \tnmark[form=enclosure, species=complement]{(6,2) .. (7,3)}{}
        \tnmark[form=enclosure, species=complement]{(6,4) .. (7,5)}{}
        \tnmark[form=enclosure, species=complement]{(6,6) .. (7,7)}{}
        \tnmark[form=label, label pos=e]{midway (2,2) and (3,2)}{$D_k$}
        \tnmark[form=label, label pos=e]{midway (4,4) and (5,4)}{$N_k$}
        \tnmark[form=bracket, label pos=e]{(2,2) .. (7,7)}{$N_{k+1}$}
    \end{tenkz}
\end{center}
